\begin{abstract}
The Theory of Structural Imagination (TSI) proposes that an internal
representation for reasoning and action should be specified by the structures it
preserves and the transformations it supports, rather than by its storage format
alone. This paper gives a consolidated formal core. A state is a typed
relational functor on a finite presented schema, a simplicial complex and metric
on a common tagged carrier, a monotone filtration, a full-support mass, and
action-conditioned tracked transitions. We define structural isomorphism and
five-layer survivor preservation, including the cross-layer compatibility axiom,
under which we prove a quantitative capacity bound on the relational defect. We
also prove equivalence and composition properties, a six-component zero criterion
including turnover, a hard-label metric--measure triangle inequality, and a
quotient-level rollout bound under explicit equivariance and Lipschitz assumptions. These are
conditional mathematical results about the specified representation class; they
do not establish that human cognition or an unconstrained neural network
implements it.
\end{abstract}
